\documentclass[10pt,a4paper]{amsart}
\usepackage[margin=19mm]{geometry}
\usepackage{amsmath,amssymb,mathtools}
\usepackage[scr=rsfs,cal=euler]{mathalfa}
\usepackage{xfrac}
\usepackage[unicode,hidelinks,bookmarks=false]{hyperref}
\newcommand{\ie}{i.e.\ }
\newcommand{\cf}{cf.\ }
\newcommand{\resp}{respectively\ }
\newcommand{\Subsection}{\subsection}
\providecommand{\nobreakdash}{\nobreak\hbox{-}}
\setlength{\emergencystretch}{2em}
\multlinegap0pt
\usepackage{xcolor}
\newtheorem{assumption}{Assumption}
\newtheorem{theorem}{Theorem}
\newcommand{\R}{\mathbb{R}}
\title{\LARGE \bf Discontinuous integro-differential equations and sliding mode control}
\author{Andrey Polyakov}
\date{Source: arXiv:2505.10116v3 (CC BY 4.0)}
\begin{document}
\maketitle

\noindent\emph{Source and licence.} \LARGE \bf Discontinuous integro-differential equations and sliding mode control. Version arXiv:2505.10116v3 (CC BY 4.0), distributed by arXiv under the Creative Commons Attribution 4.0 International licence, http://creativecommons.org/licenses/by/4.0/, as recorded in the arXiv licence metadata for this exact version and shown on the abstract page https://arxiv.org/abs/2505.10116v3. Full licence notice: \texttt{licenses/arxiv-2505.10116-CC-BY-4.0.txt}.\par\smallskip

\noindent\textit{Editorial scope.} Counted unit: the article's theorem thm:unique\_SM\_iODE (source0.tex lines 805--818). The article's complete original proof of it is delivered with it (source0.tex lines 819--865, one \texttt{proof} environment of 47 lines). No mathematics is written, repaired or re-derived: the statement and the proof are the article's own lines, in the article's own notation, and no retained line is substituted or re-flowed.  Retained uncounted as the source-local context the proof needs: lines 196--198; lines 208--210; lines 343--347; lines 351--361; lines 522--524; lines 527--529; lines 554--557; lines 572--575; lines 786--789; lines 792--799.  Nothing was omitted from any retained span, so the package carries no omission record.  No retained line contains a citation, so the wrapper carries no bibliography and no bibliography entry.  No retained line contains a figure environment, an image-inclusion command or drawing code, so no image is delivered and the package declares no asset dependency.  Editorial formatting only, in addition: an AMS \texttt{amsart} preamble that restates the article's own declarations and macros used by the retained lines, namely the environments \texttt{assumption}, \texttt{theorem} on separate counters, together with the standard packages those lines need.  The retained environments print in this document as Assumption 1, Theorem 1 in order of appearance, whereas the article prints the same items with the numbering of its own full document.  The complete native source of the article as distributed in the arXiv e-print for this exact version is bundled unchanged as \texttt{sources/178160/source0.tex}.
\begin{equation}\label{eq:ch2_generalsystem}
	\dot x(t) = f(t,x(t)), \quad f :\R \times \R^n\to  \R^n\vspace{-1mm}
\end{equation} 

\begin{equation} \label{eq:ch2_initialcondition}
	x(t^0)=x^0\in \R^n.\vspace{-1mm}
\end{equation} 

\begin{assumption}\label{as:uniq_1}\itshape
	Let  a piece-wise one-sided Lipschitz continuous function $f$ be such that 
		$f^j(t,x)\!\notin\! \mathcal{C}_{\bar G^j}(x)$ for all  $x\!\in\! \partial G^j$,
		for almost all $t\!\in\!\R$ and  for all $j\!=\!1,2,...,N$.
\end{assumption}

\begin{theorem}[\small Uniqueness of solutions in sliding mode case]\label{thm:unique_SM}\itshape
	Let\vspace{-1mm} 
		\begin{equation}
		F_{\mathcal{S}}(t,x)= K[f](t,x)\cap\mathcal{C}_{\mathcal{S}}(x), \quad t\in \R, \quad x\in \mathcal{S}.\vspace{-1mm}
		\end{equation}
 If the set 
  $F_{\mathcal{S}}$ is a singleton, i.e., $F_{\mathcal{S}}(t,x)=\{f^0(t,x)\}$ for all $x\in \mathcal{S}$, for almost all $t\in \R$, with
 for some locally {one-sided} Lipschitz continuous function $f^0$, then,  {under  Assumption \ref{as:uniq_1},}  
  the IVP  \eqref{eq:ch2_generalsystem}, \eqref{eq:ch2_initialcondition}  has a Filippov solution 
   being unique,  at least, locally in the forward time. 	
\end{theorem}

\begin{equation}\label{eq:iODE}
	\dot x(t)\!=\!f(t,x(t))\!+\!\int^t_{t^0} \!\Phi(t,\tau) \tilde f(\tau,x(\tau)) d\tau, \quad t>t^0,\vspace{-1mm}
\end{equation}

\begin{assumption}\label{as:ex_IDE}\itshape
		Let $\Phi$ be locally essentially bounded on $\Delta$ and $f,\tilde f$ satisfy Filippov Condition. 
\end{assumption}

\begin{theorem} \itshape \label{thm:uniqness_iODE} 	Let Assumption \ref{as:ex_IDE} be fulfilled and $f$, $\tilde f$ satisfy 
	Carath\'{e}\-odory  Condition. If $f$ is locally one-sided  Lipschitz  and $\tilde f$ is locally Lipschitz, then
  the IVP \eqref{eq:iODE}, \eqref{eq:ch2_initialcondition} has a  unique solution defined, at least, locally in the forward time. 
\end{theorem}	

\begin{theorem}\label{thm:sol_iODE}\itshape
If Assumption \ref{as:ex_IDE} is fulfilled, then
for any $x^0\in \R^n$ and any $t^0\in \R$,  the IVP \eqref{eq:iODE}, \eqref{eq:ch2_initialcondition} has a Filippov solution defined, at least, locally in the forward time. 
\end{theorem}

	\begin{theorem}\itshape\label{thm:uniq_one_side_IDE}
		Under Assumption \ref{as:ex_IDE},
		if $f$ is locally one-sided Lipschitz  and $\tilde f$ is locally Lipschitz  then the IVP \eqref{eq:iODE}, \eqref{eq:ch2_initialcondition} has\\ a  unique solution defined, at least, locally in the forward time. 
	\end{theorem}

	\begin{equation}\label{eq:proof_unique}\small
\begin{split}
		\!\!\!\!\!\tfrac{d \|e(t)\|^2}{dt}\!=&2e^{\!\top}\!(t)\! \left(\! f_1(t)\!-\! f_2(t)\!+\!\!\!\int^t_{t^0}\!\!\Phi(t,\tau) \left(\tilde f_1(\tau)\!-\! \tilde f_2(\tau)\right) d\tau\!\right)\\[-2mm]
	\leq& 2\ell(t)\|e(t)\|^2+2\bar \Phi \|e(t)\|\int^t_{t^0} \tilde \ell(\tau)\|e(\tau)\| d\tau\\[-2mm]
	\leq& 2\ell(t)\|e(t)\|^2\!+\!(\bar \Phi \|e(t)\|)^2\!+\!\left(\int^t_{t^0} \!\!\tilde \ell(\tau)\|e(\tau)\| d\tau\right)^2\\[-2mm]
	\leq &2\ell(t)\|e(t)\|^2+\bar \Phi^2\|e(t)\|^2+\|\tilde \ell\|_{L^2}^2\!\!\int^t_{t^0}\!\! \|e(\tau)\|^2 d\tau, 
	\end{split}\!\!\!\!\vspace{-1mm}
	\end{equation}

\begin{theorem}[\small Uniqueness of solution in sliding mode case]\label{thm:unique_SM_iODE}\hfill \newline\itshape
	Let Assumption \ref{as:ex_IDE} be fulfilled. Let    $\tilde {\mathcal{S}}\subset \R^n$ (resp., $\mathcal{S}\subset \R^n$) be a discontinuity set of the piece-wise 	(resp., one-sided) Lipschitz continuous  function $\tilde f$ (resp., $f$). Let a measurable function $(t,x,\gamma)\in \R\times \R^n\times \R^n \mapsto f^0(t,x,\gamma)\in \R^n$ be
 locally one-sided Lipschitz in $x$ and locally Lipschitz in $\gamma$.
	If \vspace{-1mm}
	\begin{itemize} 
	\item[1)]	$f$  satisfies
 Assumption \ref{as:uniq_1} and 
$\mathcal{S}\cap \tilde {\mathcal{S}}=\emptyset$;
\item[2)]  $\forall x\in \mathcal{S}$, $\exists \hat \epsilon\!>\!0$ :  $\left(K[f](t,x)\!+\!\gamma\right)\,\cap\,\mathcal{C}_{\mathcal{S}}(x)\!=\!\{f^0(t,x,\gamma)\}$ is a singleton for almost all $t\in \R$ and 
all $\gamma\in \mathcal{B}(\hat \epsilon)$;
\item[3)]  $K[f](t,x)\cap \mathcal{C}_{\tilde{\mathcal{S}}}(x)\!=\!\emptyset$ for all $x\!\in\! \tilde{\mathcal{S}}$ and almost all $t\!\in\! \R$;\vspace{-1mm}
\end{itemize}
then	the IVP  \eqref{eq:iODE}, \eqref{eq:ch2_initialcondition} has a Filippov solution being unique,  at least, locally in the forward time. 	
\end{theorem}

\begin{proof}
	Let us study the case $x^0\in \mathcal{S}$. 	This means that $x^0\in \partial G^j$ for all $j\in \mathcal{N}(x^0)$, where, as before, $\mathcal{N}:\mathcal{S}\to 2^{\{1,2,...,N\}}$ is the indicator function, i.e., $j\in \mathcal{N}(x_0)$ if $x_0\in \partial G^{j}$.
	Since the function $f$ is piece-wise continuous then  there exists $\bar \epsilon>0$ 
	such that $\mathcal{N}(x^0+\delta)\subset \mathcal{N}(x^0)$ for $x_0+\delta\in S: \|\delta\|\leq \bar \epsilon$. 
	
	
	
	By 
	Assumption \ref{as:uniq_1}, the vector $f^j(t,x^0)$ belongs to the set $\R^n\backslash \mathcal{C}_{\bar  G^{j}}(x^0)$ 
	for all $t\in \R$.
	Since $\mathcal{C}_{\bar G^{j}}(x^0)$ is a positive closed cone then its complement $\R^n\backslash \mathcal{C}_{\bar G^{j}}(x^0)$ is a positive open cone.
	Hence, taking into account,
	$f^j\in C(\R\times \R^n;\R^n)$  we conclude that there exists $\epsilon^j>0$ such that \vspace{-1mm}
	$$
    f^j(t^0+s,x^0+\mathcal{B}(\epsilon^j))+\mathcal{B}(\epsilon^j) \subset  \left(\R^n\backslash {\mathcal{C}}_{\bar G^{j}}(x^0)\right), 	\; \forall s\in (0,\epsilon^j),\vspace{-1mm}
	$$
 where $\mathcal{B}(\epsilon^j)\subset \R^n$ is an open  ball of the radius $\epsilon^j>0$.
	
	Let $\epsilon =\min\left\{\bar \epsilon,\hat \epsilon, \min_{j\in \mathcal{N}(x^0)}\{\epsilon^j\}\right\}$.
	Repeating the proof of Theorem \ref{thm:unique_SM},  we conclude  that any solution 
	$x^{\epsilon}$
	of the IVP\vspace{-1mm}
	\begin{equation}\label{eq:ODI_temp_unique_SM}
		\dot x^{\epsilon}(t)\in K[f](t,x^{\epsilon}(t))+\mathcal{B}(\epsilon), \quad x^{\epsilon}(t^0)=x^0\in \mathcal{S}\vspace{-1mm}
	\end{equation}
	belongs to $\mathcal{S}$	as long as $x^{\epsilon}(t)\in x^0+\mathcal{B}(\epsilon)$. 
	For any continuous function $t\in [t^0,t^0+\epsilon]\to x(t)\in \R^n$, the set-valued function \vspace{-1mm}
	$$
	t\in [t^0,t^0+\epsilon]\mapsto \tilde Q(t):=\int^t_{t^0} \Phi(t,\tau) K[\tilde f](\tau,x(\tau)) d\tau\subset \R^n\vspace{-1mm}
	$$ is bounded as follows (see, the proof of Theorem \ref{thm:sol_iODE})\vspace{-1mm}
	\[
	\sup_{y\in \tilde Q(t)} \|y\|\leq \overline{\Phi}\int^t_{t^0}m(\tau) +\eta\, d\tau\vspace{-1mm}
	\]
	for some $\eta>0, \overline{\Phi}>0$ and some integrable function $m:[t^0,t^0+\epsilon]\to \R$. This means that $\exists t'\in (0,\epsilon)$ such that \vspace{-1mm}
	\[
	\sup_{y\in \tilde Q(t)} \|y\|\leq\epsilon, \quad \forall t\in [t^0,t^0+t'].\vspace{-1mm}
	\]
	Any solution the system \eqref{eq:iODE} with $x(t^0)=x^0\in \mathcal{S}$ is a solution of \eqref{eq:ODI_temp_unique_SM} for  $t\in[t^0,t^0+t']$. Since $\mathcal{S}\cap \tilde {\mathcal{S}}=\emptyset$  and  $\tilde f$ satisfies Lipschitz condition on any compact from $\mathcal{S}$  then  the IDE\vspace{-1mm}
	\begin{equation*}
\dot x(t)\!=\!f^0\left(t,x(t),\gamma(t)\right)+\gamma(t), \quad \gamma(t)\!=\!\int^t_{t^0}\!\!\!\Phi(t,\tau) \tilde f(\tau,x(\tau)) d\tau\vspace{-1mm}
	\end{equation*}
   defines the motion of the system on the sliding surface. 
  {For the obtained system, repeating the proof of Theorem  \ref{thm:uniq_one_side_IDE}, we derive the estimate similar to \eqref{eq:proof_unique}, which insures a local in-time uniqueness of solutions.  	For $x^0\in  G^j\backslash \tilde{\mathcal{S}}$, the local-in-time  uniqueness of solution  follows from  Theorem \ref{thm:uniqness_iODE}.  Finally, since $K[f](t,x)$ and $\mathcal{C}_{\mathcal{S}}(x)$ are compact sets  for all $x$ and for almost all $t$, then the condition 3) implies that 
  	$\forall x\in \tilde {\mathcal{S}}$, $\exists \hat \epsilon\!>\!0$ such that  $\left(K[f](t,x)+\gamma \right)\cap \mathcal{C}_{\tilde{\mathcal{S}}}(x)=\emptyset$  for almost all $t\in \R$ and 
  	all $\gamma\in \mathcal{B}(\hat \epsilon)$.
  	So, any solution  initiated on $\tilde S$ immediately enters into  domain $\R^n\backslash\mathcal{S}$, i.e., the estimate \eqref{eq:proof_unique} holds almost everywhere and solution is unique, at least, locally in time.}
\end{proof}

\end{document}
